\chapter{精选习题详解}

下面的计算各自给出完整题设，不依赖已经失效的“原专题”编号。它们选取正文中容易把定义域、归一化或条件漏掉的练习。读者可以先独立求解，再逐步核对。

\section{群作用、函数空间与边界条件}

\subsection*{正方形顶点的轨道与稳定子}
设循环群 \(C_4=\{e,r,r^2,r^3\}\) 以绕中心旋转作用于正方形的四个顶点，且 \(r^4=e\)。固定顶点 \(v\)，求轨道 \(C_4v\) 与稳定子 \(C_{4,v}=\{g\in C_4:gv=v\}\)。

四次旋转分别把 \(v\) 送到四个不同顶点，所以 \(|C_4v|=4\)。只有恒等元固定 \(v\)，故 \(C_{4,v}=\{e\}\)。轨道--稳定子计数式在这里成为 \(4=4/1\)。若把作用对象换成中心点，稳定子就变成整个 \(C_4\)，轨道则只有一个点；它们依赖所选的点，而不只依赖群。

\subsection*{\(L^2\) 函数不能按点取值}
在 \(L^2(0,1)\) 中令 \(f_n(x)=\max(1-nx,0)\)。计算 \(f_n(0)\) 与范数，并判断“取 \(x=0\) 处的值”能否成为 \(L^2\) 上的连续线性泛函。

每个 \(f_n\) 都连续且 \(f_n(0)=1\)，但
\[
\|f_n\|_2^2=\int_0^{1/n}(1-nx)^2\dd x
=\frac1{3n}\longrightarrow0.
\]
如果存在连续线性泛函 \(T:L^2(0,1)\to\mathbb C\)，并且对连续函数总有 \(T(f)=f(0)\)，连续性将推出 \(T(f_n)\to T(0)=0\)，与 \(T(f_n)=1\) 矛盾。更根本地说，\(L^2\) 的元素是几乎处处相等的函数的等价类，单点值并不是类的性质。要讨论边界值，必须先给函数增加足够的正则性，例如取一维 \(H^1(0,1)\) 的迹。

\subsection*{Robin 边界条件为何给自伴算子}
在 \(L^2(0,1)\) 上取 \(S_Rf=-f''\)，定义域为
\[
D(S_R)=\{f\in H^2(0,1):f'(0)=\alpha f(0),\ f'(1)=-\beta f(1)\},
\qquad \alpha,\beta\in\mathbb R.
\]
证明它自伴。这里内积对第二个变量线性。

两次分部积分给出
\[
\langle S_Rf,g\rangle-\langle f,S_Rg\rangle
=\bigl[\overline{f}g'-\overline{f'}g\bigr]_0^1.
\]
若 \(f,g\in D(S_R)\)，则每个端点处的两项因 \(\alpha,\beta\) 为实数而相消，所以 \(S_R\) 对称。反过来，若 \(g\in D(S_R^*)\)，先用紧支撑的 \(f\) 检验，可得分布意义下 \(g''\in L^2\)，从而 \(g\in H^2(0,1)\)。再让 \(f(0),f(1)\) 任意变化；满足给定 Robin 条件的三次 Hermite 插值多项式足以实现这样的端点值。边界型因而变成
\[
\overline{f(1)}[g'(1)+\beta g(1)]
-\overline{f(0)}[g'(0)-\alpha g(0)]=0.
\]
两个系数必须分别为零，即 \(g\in D(S_R)\)。于是 \(D(S_R^*)=D(S_R)\)，且 \(S_R^*=S_R\)。只写出边界型为零只能证明对称；还要检查伴随定义域，才能证明自伴。

\subsection*{Neumann 算子的零模与可解条件}
求解 \(-u''=f\) 于 \((0,1)\)，并施加 \(u'(0)=u'(1)=0\)。说明存在解的条件及解的唯一性。

对方程积分，利用两个边界条件得到 \(\int_0^1f(x)\dd x=0\)。反之，若 \(f\in L^2(0,1)\) 的积分为零，令
\[
u'(x)=-\int_0^x f(s)\dd s,\qquad
u(x)=C-\int_0^x(x-s)f(s)\dd s.
\]
它满足方程和两个 Neumann 条件，因此条件也充分。任意常数 \(C\) 都给出一个解；若两个解相减，其二阶导数为零且两端导数为零，所以差必为常数。零模不是计算瑕疵，而是算子不能在整个 \(L^2\) 上求逆的原因。

\subsection*{半直线动量与二阶算子的亏量}
在 \(L^2(0,\infty)\) 上分别从 \(C_c^\infty(0,\infty)\) 出发考虑 \(P=-\ii\partial_x\) 与 \(A=-\partial_x^2\)。求它们的亏量维数，并判断是否可能存在自伴扩张。亏量空间定义为 \(\ker(T^*\mp\ii)\)。

对动量，\((P^*-\ii)f=0\) 化为 \(f'=-f\)，其解 \(e^{-x}\) 平方可积；\((P^*+\ii)f=0\) 化为 \(f'=f\)，其非零解 \(e^x\) 不平方可积。因此 \((n_+,n_-)=(1,0)\)，两数不等，故不存在自伴扩张。对二阶算子，\((A^*-\ii)f=0\) 与 \((A^*+\ii)f=0\) 分别给 \(f''=-\ii f\) 与 \(f''=\ii f\)。每个方程的两个指数解中恰有一个随 \(x\to\infty\) 衰减，所以 \((n_+,n_-)=(1,1)\)。二阶算子因此允许自伴扩张，例如端点的 Dirichlet 条件或实 Robin 条件。两个微分表达式都能在紧支撑函数上做形式分部积分，但它们在半直线上的自伴性结论完全不同。

\section{源、传播子与不稳定的逆}

\subsection*{阶跃函数的分布导数}
设 \(H\) 为 Heaviside 函数，\(f(x)=H(x)-H(x-1)\)。求 \(f'\) 作为分布的表达式。

对任意紧支撑光滑测试函数 \(\varphi\)，分布导数的定义给
\[
\langle f',\varphi\rangle=-\langle f,\varphi'\rangle
=-\int_0^1\varphi'(x)\dd x
=\varphi(0)-\varphi(1).
\]
因此 \(f'=\delta_0-\delta_1\)。端点处如何人为规定 \(H(0)\) 不影响答案。若 Fourier 约定为 \(\widehat f(k)=\int e^{-\ii kx}f(x)\dd x\)，则 \(\widehat{f'}(k)=\ii k\widehat f(k)=1-e^{-\ii k}\)，与两个点源的变换相同。

\subsection*{二维 Poisson 方程的点源常数}
在 \(\mathbb R^2\) 令 \(E(x)=-(2\pi)^{-1}\log|x|\)。验证它在原点以外满足 \(\Delta E=0\)，并求围绕原点的小圆的外法向通量。

对径向函数，\(\Delta E=E''(r)+r^{-1}E'(r)\)。由于 \(E'(r)=-(2\pi r)^{-1}\)、\(E''(r)=(2\pi r^2)^{-1}\)，两项在 \(r>0\) 处抵消。半径为 \(\varepsilon\) 的圆长为 \(2\pi\varepsilon\)，故
\[
-\int_{|x|=\varepsilon}\partial_rE\dd s
=\frac1{2\pi\varepsilon}(2\pi\varepsilon)=1.
\]
把 \(\mathbb R^2\) 切去小圆，对测试函数分部积分，再令 \(\varepsilon\downarrow0\)，这项通量留下 \(\varphi(0)\)；其余小圆边界项随 \(\varepsilon|\log\varepsilon|\) 消失。故严格的分布等式是 \(-\Delta E=\delta_0\)，而不只是原点外的调和性。

\subsection*{有限秩 Fredholm 方程}
令 \(H\) 为 Hilbert 空间，内积对第二变量线性。取 \(a_j,b_j\in H\)，并定义 \(Kx=\sum_{j=1}^m b_j\langle a_j,x\rangle\)。把 \((I-K)u=f\) 化为有限维线性方程，并解释奇异矩阵意味着什么。

记 \(c_i=\langle a_i,u\rangle\)。原方程等价于
\[
u=f+\sum_{j=1}^m b_jc_j,\qquad
c_i=\langle a_i,f\rangle+\sum_{j=1}^m\langle a_i,b_j\rangle c_j.
\]
写 \(M_{ij}=\langle a_i,b_j\rangle\)、\(d_i=\langle a_i,f\rangle\)，便得到 \((I-M)c=d\)。只要这个线性方程可解，代回的 \(u\) 即为原方程的解；反过来原方程的任何解也给出这样的 \(c\)。若 \(I-M\) 不可逆，齐次方程具有非零 \(c\)，相应 \(u=\sum b_jc_j\) 不可能为零，因为 \(c_i=\langle a_i,u\rangle\)。非齐次问题此时还需 \(d\) 与左零空间正交，不能仅凭“矩阵奇异”断言无解。

\subsection*{Tikhonov 滤波怎样抑制高频}
在 \(\ell^2\) 上取 \(Ke_n=n^{-1}e_n\)。给定带噪数据 \(y^\delta\)，令 \(x_\alpha=(K^*K+\alpha I)^{-1}K^*y^\delta\)，其中 \(\alpha>0\)。求每个分量，比较它与精确逆的噪声放大。

因为 \(K^*=K\) 且 \(Ke_n=n^{-1}e_n\)，故
\[
(x_\alpha)_n=\frac{n}{1+\alpha n^2}\,y_n^\delta.
\]
无正则化时形式逆把第 \(n\) 个噪声分量乘 \(n\)；正则化因子在 \(n\gg\alpha^{-1/2}\) 时却近似 \(1/(\alpha n)\)。对无噪数据 \(y=Kx\)，重建的第 \(n\) 个分量是 \(x_n/(1+\alpha n^2)\)，偏差为 \(-\alpha n^2x_n/(1+\alpha n^2)\)。因此较大的 \(\alpha\) 增强稳定性，也更早压低真实高频。

\subsection*{Burgers 激波的熵产生}
对 \(u_t+(u^2/2)_x=0\)，令左右常态为 \(u_L=2,u_R=0\)。求不连续面 \(x=st\) 的速度，并用凸熵 \(\eta(u)=u^2/2\) 检验它。相应熵通量由 \(q'(u)=\eta'(u)f'(u)\) 定义。

Rankine--Hugoniot 条件给 \(s=(f(2)-f(0))/(2-0)=1\)。因为 \(q'(u)=u^2\)，可取 \(q(u)=u^3/3\)。对跳跃量约定 \([h]=h_R-h_L\)，分布 \(\partial_t\eta+\partial_xq\) 在 \(x=st\) 上的系数是 \([q]-s[\eta]\)，此处为
\[
\left(0-\frac83\right)-1(0-2)=-\frac23<0.
\]
因此它满足熵不等式 \(\partial_t\eta+\partial_xq\le0\)。把左右状态倒过来时，跳跃条件仍给 \(s=1\)，但同一系数变成 \(+2/3\)；这个不连续弱解违反熵条件，物理解应改为特征张开的稀疏扇。

\section{几何对象不能只凭坐标判断}

\subsection*{圆柱面的切向量}
圆柱面由 \(F(x,y,z)=x^2+y^2-1=0\) 给出。取局部参数化 \(X(\theta,z)=(\cos\theta,\sin\theta,z)\)，求两个坐标切向量并验证它们属于约束的核。

直接求导得
\[
X_\theta=(-\sin\theta,\cos\theta,0),\qquad X_z=(0,0,1).
\]
由于 \(dF_p(v)=2xv_x+2yv_y\)，代入 \(p=X(\theta,z)\) 可得 \(dF_p(X_\theta)=dF_p(X_z)=0\)。两向量线性无关，而正则约束 \(dF_p\ne0\) 的核维数为二，故它们张成整个切空间。\(\theta\) 只能在局部区域取连续实值：绕圆柱一圈会让角度增加 \(2\pi\)，所以单个实值角坐标不能覆盖整个圆柱。

\subsection*{球面曲率在整体缩放下怎样变化}
单位球面的度量是 \(g_1=\dd\vartheta^2+\sin^2\vartheta\,\dd\varphi^2\)。令 \(g_a=a^2g_1\)，其中 \(a>0\)。求 Gauss 曲率。

取正交余标架 \(\theta^1=a\dd\vartheta\)、\(\theta^2=a\sin\vartheta\,\dd\varphi\)。Cartan 第一结构方程 \(\dd\theta^i+\omega^i{}_j\wedge\theta^j=0\) 给 \(\omega^1{}_2=-\cos\vartheta\,\dd\varphi\)（符号与取向约定一致）。因二维反对称联络没有 \(\omega\wedge\omega\) 项，
\[
\Omega^1{}_2=\dd\omega^1{}_2
=\sin\vartheta\,\dd\vartheta\wedge\dd\varphi
=\frac1{a^2}\theta^1\wedge\theta^2.
\]
按照 \(\Omega^1{}_2=K\theta^1\wedge\theta^2\) 的定义，\(K=1/a^2\)。连接形式没有随常数 \(a\) 改变，但单位面积变为原来的 \(a^2\) 倍；这正是曲率缩放的来源。

\subsection*{不可积的平面分布}
在 \(\mathbb R^3\) 上令 \(\alpha=\dd z-x\dd y\)，并令二维分布 \(D=\ker\alpha\)。判断它能否在每一点附近成为某族曲面的切平面。

向量场 \(X=\partial_x\)、\(Y=\partial_y+x\partial_z\) 都满足 \(\alpha(X)=\alpha(Y)=0\)，且处处线性无关，因而张成 \(D\)。由交换子的定义，
\[
[X,Y]f=X(Yf)-Y(Xf)=\partial_z f,\qquad [X,Y]=\partial_z.
\]
但 \(\alpha(\partial_z)=1\)，所以交换子不属于 \(D\)。若 \(D\) 是曲面的切平面，沿曲面切向量场的交换子仍必须切于曲面，矛盾。等价地，\(\dd\alpha=-\dd x\wedge\dd y\) 且 \(\alpha\wedge\dd\alpha=-\dd z\wedge\dd x\wedge\dd y\ne0\)，违背 Frobenius 条件。

\subsection*{平移不变性与应力能量守恒}
在平直时空取实标量场的密度 \(\mathcal L=\tfrac12\partial_\mu\phi\,\partial^\mu\phi-\tfrac12\kappa^2\phi^2\)，其中 \(\kappa>0\) 为常数。定义 Euler--Lagrange 表达式 \(E(\phi)=\partial\mathcal L/\partial\phi-\partial_\mu[\partial\mathcal L/\partial(\partial_\mu\phi)]\)。验证正则张量 \(T^\mu{}_\nu=\partial^\mu\phi\,\partial_\nu\phi-\delta^\mu{}_\nu\mathcal L\) 的散度恒等式。

先有 \(E(\phi)=-(\Box+\kappa^2)\phi\)。对 \(T^\mu{}_\nu\) 求导并用乘积法则，
\[
\partial_\mu T^\mu{}_\nu
=(\Box\phi)\partial_\nu\phi
+\partial^\mu\phi\,\partial_\mu\partial_\nu\phi
-\partial_\nu\mathcal L.
\]
第二项等于 \(\partial_\nu[\tfrac12\partial_\rho\phi\,\partial^\rho\phi]\)，因此与 \(\partial_\nu\mathcal L\) 中的动能导数抵消，留下
\[
\partial_\mu T^\mu{}_\nu
=(\Box\phi+\kappa^2\phi)\partial_\nu\phi
=-E(\phi)\partial_\nu\phi.
\]
这在任意光滑场上都是恒等式；只有再施加运动方程 \(E(\phi)=0\)，才得到守恒律。

\section{随机变量与可识别的参数}

\subsection*{收敛方式的区别}
设 \(P(A_n)=n^{-2}\)，并取 \(X_n=n\mathbf1_{A_n}\)。判断 \(X_n\) 是否在 \(L^1\)、\(L^2\) 以及概率意义下收敛到零。

由定义，\(\mathbb E|X_n|=nP(A_n)=1/n\to0\)，所以 \(X_n\to0\) 于 \(L^1\)。另一方面，\(\mathbb E|X_n|^2=n^2P(A_n)=1\)，故不在 \(L^2\) 收敛。对任意固定 \(\varepsilon>0\)，当 \(n>\varepsilon\) 时 \(P(|X_n|>\varepsilon)=P(A_n)=n^{-2}\to0\)，因此仍依概率收敛。事件之间是否独立对这三个结论没有影响。

\subsection*{Bernoulli 样本总数的完备性}
设 \(X_1,\ldots,X_n\) 独立且服从 Bernoulli\((p)\)，\(0<p<1\)，并记 \(T=\sum_iX_i\)。证明 \(T\) 完备：若函数 \(a:\{0,\ldots,n\}\to\mathbb R\) 对所有 \(p\) 满足 \(\mathbb E_p a(T)=0\)，则 \(a(t)=0\) 对每个 \(t\) 成立。

二项分布给出
\[
0=\sum_{t=0}^n a(t){n\choose t}p^t(1-p)^{n-t}.
\]
除以正数 \((1-p)^n\)，再令 \(z=p/(1-p)\in(0,\infty)\)，得到多项式恒等式
\[
\sum_{t=0}^n a(t){n\choose t}z^t=0\qquad(z>0).
\]
一个多项式在无穷多个点为零，只能每个系数都为零；又因组合数非零，故 \(a(t)=0\)。证明依赖参数 \(p\) 在整个开区间变化；若只固定一个 \(p\)，零均值函数当然不必处处为零。

\subsection*{参数依赖支撑集为何破坏常规 Fisher 等式}
设 \(X\sim\operatorname{Uniform}(0,\theta)\)，\(\theta>0\)。若只在内部 \(0<x<\theta\) 对对数密度求导，会得到什么“得分”？它为何不能直接套用常规 Cramér--Rao 证明？

内部密度为 \(1/\theta\)，故形式得分 \(\partial_\theta\log f_\theta(x)=-1/\theta\)。若误用“得分均值为零”，便得到 \(-1/\theta=0\) 的荒谬等式。真正的归一化恒等式是
\[
1=\int_0^\theta\frac1\theta\dd x,\qquad
0=\frac{\dd}{\dd\theta}\int_0^\theta\frac1\theta\dd x
=\int_0^\theta\left(-\frac1{\theta^2}\right)\dd x+\frac1\theta.
\]
最后一项来自随参数移动的端点，正好抵消内部积分。常规证明把导数移入一个固定支撑集上的积分，因此在此模型中不合法；不能把内部形式得分的平方均值 \(1/\theta^2\) 直接当作适用常规下界的 Fisher 信息。

\subsection*{正态线性模型的残差自由度}
设 \(Y\sim N(X\beta,\sigma^2I_n)\)，其中 \(X\) 是 \(n\times p\) 的满列秩设计矩阵，\(p<n\)。令 \(H=X(X^\top X)^{-1}X^\top\)。证明 \(H\) 是正交投影，并求残差空间的维数。

矩阵 \(X^\top X\) 对称可逆，因此 \(H^\top=H\)。直接相乘又得
\[
H^2=X(X^\top X)^{-1}(X^\top X)(X^\top X)^{-1}X^\top=H.
\]
对称且幂等的矩阵是正交投影，其像空间恰为 \(\operatorname{col}(X)\)，维数为 \(p\)。残差 \((I-H)Y\) 落在其正交补内，故残差自由度为 \(\operatorname{rank}(I-H)=n-p\)。若进一步使用正态性，\(HY\) 与 \((I-H)Y\) 的协方差为 \(\sigma^2H(I-H)=0\)，联合正态遂使二者独立；没有正态性时“零协方差推出独立”不成立。
